\scanpage{16}\mspage{15}
Supports et cosupports.

$F$ et $G$ de $\unsure{\mathcal F|\mathscr F}$ sont \uline{disjoints} \inserted{(au sens de $\leq$)} sss $\lvert F\rvert\cap\lvert G\rvert=\emptyset$. Donc les $F$ disjoints d'une famille de figures $G_i$, sont celles qui \deleted{sont} satisfont $\lvert F\rvert\subset L'$, où $L'=L\setminus\bigcup_i\lvert G_i\rvert$. Donc, on voit que les supports correspondent aux parties quelconques $L'$ de $L$, la \unsure{relation|notion} de complémentarité est la notion ensembliste.

\deleted{Le $\mathcal F(L'$} L'ens des figures \inserted{de $L'$} de support $\subset L'$ \deleted{est} s'identifie aux figures de $L'$, i.e. on trouve les $\mathcal F(L')\subset\mathcal F(L)$ comme sous-ens. saturés (au sens de $\leq$ et de la disjonction \ldots).
% Paper watermark; sideways and faint; partial impressions, not manuscript prose.
\margin{\unsure{RONG}\quad\unsure{EXTRA STR}}
